\documentclass[level=2]{kidscode}
% Let lines with long code words stretch a little instead of running into the margin.
\setlength{\emergencystretch}{3em}
\newcommand{\bx}{$\square$}
\begin{document}
\kctitle{Final Project Brief}{Arcade Game or Art Studio: choose, plan, build, test and show (weeks 9--12)}

\noindent\textbf{Your mission:} you have learned functions, lists, dictionaries, loops, files, tests and games with keys and
timers. Now you put them together in a program that is completely \textbf{yours}. In week 12
your family will come to our \textbf{Family Arcade} to play your game or make art with your
studio, and you will explain how it works. You can work alone or with a partner.

\section{Step 1: Choose your option}

\begin{center}
%% alt: Two option cards side by side. Option A, arcade game: a black game screen showing the score and lives, a small pointy star, a big round meteor and a basket at the bottom that you move with the keys. Option B, art studio: a flower made of eight green petals on a white canvas, a saved-design note that reads petals 8, size 120, color green, and the keys n for new, s for save and l for load.
\begin{tikzpicture}[font=\small,scale=0.92]
  % Option A: arcade game
  \draw[rounded corners=4pt, thick, fill=kcsky!10] (0,0) rectangle (7.4,4.6);
  \node[anchor=north west, font=\bfseries\large] at (0.2,4.5) {Option A: Arcade game};
  \fill[kcblack] (0.5,0.5) rectangle (6.9,3.7);
  \fill[kcsky, draw=white] (2.6,0.8) rectangle (4.0,1.05);
  \node[text=white, font=\scriptsize] at (3.3,0.62) {basket (keys)};
  \node[star, star points=5, star point ratio=2.3, fill=kcyellow, minimum size=0.45cm, inner sep=0pt] (s) at (2.0,2.8) {};
  \node[text=white, font=\scriptsize, anchor=west] at (2.3,2.8) {star: small, pointy};
  \fill[kcvermillion] (5.3,2.1) circle (0.32);
  \node[text=white, font=\scriptsize, anchor=north] at (5.3,1.75) {meteor: big, round};
  \node[text=kcyellow, font=\scriptsize\bfseries, anchor=north west] at (0.6,3.65) {Score: 7   Lives: 2};
  \node[anchor=north west, text width=7cm, font=\scriptsize] at (0.2,0.45) {};
  % Option B: art studio
  \begin{scope}[xshift=8cm]
  \draw[rounded corners=4pt, thick, fill=kcpurple!10] (0,0) rectangle (7.4,4.6);
  \node[anchor=north west, font=\bfseries\large] at (0.2,4.5) {Option B: Art studio};
  \fill[white] (0.5,0.5) rectangle (4.6,3.7);
  \foreach \a in {0,45,...,315}
    \draw[kcgreen, thick, rotate around={\a:(2.55,2.1)}] (2.55,2.1) .. controls (3.0,2.5) and (3.4,2.5) .. (3.6,2.1) .. controls (3.4,1.7) and (3.0,1.7) .. (2.55,2.1);
  \draw[thick, fill=white] (4.8,1.5) -- (6.7,1.5) -- (7.2,2.0) -- (7.2,3.5) -- (4.8,3.5) -- cycle;
  \node[font=\scriptsize\ttfamily, align=left] at (6.0,2.55) {petals=8\\size=120\\color=green};
  \node[font=\scriptsize, anchor=north] at (6.0,1.45) {saved design};
  \node[font=\scriptsize, anchor=north west] at (0.5,0.45) {keys: n = new, s = save, l = load};
  \end{scope}
\end{tikzpicture}
\end{center}

\begin{tabularx}{\linewidth}{@{}>{\bfseries}p{3.2cm}XX@{}}
\toprule
 & Option A: Arcade game & Option B: Art studio \\ \midrule
What it is & A turtle game with \textbf{keys}, a \textbf{timer} and a \textbf{score}. & A program that \textbf{makes designs} (often random) and can \textbf{save} and \textbf{load} them. \\
You will love it if\ldots & you like playing and inventing games. & you like patterns, colors and surprises. \\
Starter file & \texttt{starter/\allowbreak game\_template.py} & \texttt{starter/\allowbreak art\_studio\_template.py} \\
The file saves\ldots & high scores & designs (one setting per line) \\
\bottomrule
\end{tabularx}

\section{Step 2: Know the must-haves}
Whatever you choose, your finished program must have all six. Check them off as you go.
\begin{checklist}
  \item \textbf{Functions} with parameters that \textbf{return} a value (for example \texttt{move(x, step)} or \texttt{is\_touching(a, b)}).
  \item A \textbf{list} (enemies, coins, colors, a gallery of saved designs\ldots).
  \item A \textbf{dictionary} (settings, a player, a design\ldots).
  \item A \textbf{while loop} or a \textbf{game loop with} \py|ontimer|.
  \item A \textbf{file}: high scores or saved designs, next to your program, created if it is missing.
  \item At least \textbf{3 pytest tests} of your own functions, in a test file (\texttt{my\_tests.py} from the
        template, or your own \texttt{test\_mygame.py}) that you run with \menu{Run}.
\end{checklist}

\begin{tip}[The secret of easy tests]
Put the \emph{rules} of your program (scores, moving, touching, saving) in small functions that do
\textbf{not} use turtle. They take values in and return an answer, so they are easy to test.
Keep the drawing in other functions.
\end{tip}

\section{Step 3: Get an idea}
Pick one from the menu, change it, or invent your own.

{\small
\renewcommand{\arraystretch}{1.3}
\begin{longtable}{@{}>{\bfseries\raggedright\arraybackslash}p{2.8cm}>{\raggedright\arraybackslash}p{4.4cm}>{\raggedright\arraybackslash}p{3.6cm}>{\raggedright\arraybackslash}p{3.6cm}@{}}
\toprule
Idea & What happens & A list of\ldots & The dictionary / file \\ \midrule\endhead
\multicolumn{4}{@{}l}{\textbf{Arcade games}} \\ \midrule
Space Snack & Fly a ship left and right, catch snacks, dodge asteroids. & falling snacks and asteroids & settings; high scores \\
Dragon Feeder & Move a dragon to eat apples before the timer ends. Golden apples are worth more. & apples on the screen & points for each kind; high scores \\
Rainy Day & Move an umbrella to keep a cat dry. Rain falls faster over time. & raindrops & speed levels; best time \\
Bug Squasher & Bugs appear in random places; move a net onto them with the arrow keys. & bugs & bug types and points; high scores \\
Treasure Diver & A diver collects pearls while sharks swim side to side. & sharks and pearls & settings; top 5 scores \\ \midrule
\multicolumn{4}{@{}l}{\textbf{Art studios}} \\ \midrule
Kaleidoscope Lab & Random mirrored patterns; press keys to change the number of mirrors. & colors, shapes & one design; saved designs file \\
Pixel Pet Maker & Draw a pet on a grid of squares from rows of letters. & rows of the grid & letter $\to$ color; saved pets \\
Spiro Studio & Spirograph-style loops; keys change size and turn. & a gallery of designs & one design; saved gallery \\
Quilt Maker & A quilt of squares, each with a random pattern. & patterns & pattern settings; saved quilts \\
Firework Show & Click or press keys to launch fireworks; save your show and replay it. & fireworks in the show & one firework; saved show \\
\bottomrule
\end{longtable}}

\section{Step 4: Milestones}
Big projects are built in small steps. Show your teacher each milestone.

{\renewcommand{\arraystretch}{1.35}
\begin{longtable}{@{}>{\bfseries\raggedright\arraybackslash}p{3.0cm}>{\raggedright\arraybackslash}p{2.6cm}>{\raggedright\arraybackslash}p{8.0cm}>{\centering\arraybackslash}p{0.8cm}@{}}
\toprule
Milestone & When & Done when\ldots & \faCheck \\ \midrule\endhead
M1 Pitch and plan & week 9 (pitch), start of week 10 (plan) & Your pitch page, your pseudocode or flowchart, and (in a pair) your task split are filled in. & \bx \\ \midrule
M2 Core loop works & week 10 & The main thing happens on screen: the player moves and things fall, or the studio draws a design when you press a key. & \bx \\ \midrule
M3 Data and tests & week 11, first half & Your list, dictionary and file work. Your test file has 3 tests that pass. & \bx \\ \midrule
M4 Polish & week 11, second half & Someone playtested it, you fixed one thing they found, and you did the accessibility check. & \bx \\ \midrule
Showcase & week 12 & You showed it at the Family Arcade and explained one part. & \bx \\
\bottomrule
\end{longtable}}

%% ---------------------------------------------------------------------
\clearpage
\section{Planning page 1: My pitch (M1)}
\namedateline
Project name: \rule{9cm}{0.4pt}\qquad Option: A \bx\quad B \bx

\noindent In one sentence, my project is\ldots
\answerlines{2}
\noindent How do you play it or use it? (keys, goal, how it ends)
\answerlines{2}
\noindent My \textbf{list} will hold\ldots \answerlines{1}
\noindent My \textbf{dictionary} will hold\ldots \answerlines{1}
\noindent My \textbf{file} will save\ldots \answerlines{1}
\noindent My \textbf{loop} (\texttt{while} or \texttt{ontimer}) will\ldots \answerlines{1}
\noindent One thing that will make it special:
\answerlines{1}

\section{Planning page 2: Pseudocode or flowchart (M1)}
Plan the steps in plain words or boxes and arrows before you write Python. Here is the plan for a
catching game:

\begin{center}
%% alt: Flowchart of the game loop. Start with score 0 and lives 3. Move everything down. Is anything touching the basket? If yes, a star adds 1 to the score and a meteor takes 1 from the lives. Then ask: are the lives 0? If yes, the game is over and the score is saved in the file. If no, wait 40 milliseconds with ontimer and go back to moving everything down.
\begin{tikzpicture}[
  node distance=0.45cm and 1.3cm, font=\small,
  box/.style={draw, thick, rounded corners=3pt, fill=kcsky!15, align=center, minimum height=0.8cm, text width=3.9cm},
  ask/.style={draw, thick, diamond, aspect=2.4, fill=kcyellow!35, align=center, inner sep=1pt, text width=2.2cm},
  arrow/.style={-{Stealth[length=2.5mm]}, thick}]
  \node[box, fill=kcgreen!20] (start) {Start: score 0, lives 3};
  \node[box, below=of start] (move) {Move everything down};
  \node[ask, below=of move] (touch) {Touching the basket?};
  \node[box, right=of touch, text width=3.2cm] (score) {Star: score + 1\\Meteor: lives $-$ 1};
  \node[ask, below=of touch] (over) {Lives = 0?};
  \node[box, right=of over, text width=3.2cm, fill=kcorange!25] (end) {Game over: save the score in the file};
  \node[box, left=of touch, text width=2.4cm] (wait) {Wait 40 ms\\(\texttt{ontimer})};
  \draw[arrow] (start) -- (move);
  \draw[arrow] (move) -- (touch);
  \draw[arrow] (touch) -- node[above]{yes} (score);
  \draw[thick] (score.south) |- ($(touch.south)!0.5!(over.north)$);
  \draw[arrow] (touch) -- node[left]{no} (over);
  \draw[arrow] (over) -- node[above]{yes} (end);
  \draw[arrow] (over) -| node[pos=0.25,above]{no} (wait);
  \draw[arrow] (wait) |- (move);
\end{tikzpicture}
\end{center}

\noindent My plan (pseudocode or flowchart):
\answerspace{7.5cm}

\section{Planning page 3: Task split for pairs (M1)}
Split the work fairly: each person gets interesting jobs, and each person writes at least one
function \textbf{and} one test. Swap driver and navigator every 8 minutes as usual.

{\renewcommand{\arraystretch}{1.6}
\begin{tabularx}{\linewidth}{@{}>{\raggedright\arraybackslash}p{6.4cm}|X|>{\centering\arraybackslash}p{1.2cm}@{}}
\toprule
\textbf{Task} & \textbf{Who} & \textbf{Done} \\ \midrule
Rules functions (move, touching, score, \ldots) & & \bx \\ \midrule
Drawing: shapes, colors, text on screen & & \bx \\ \midrule
Keys and timer (\texttt{onkey}, \texttt{ontimer}) or key menu & & \bx \\ \midrule
The list and the dictionary & & \bx \\ \midrule
Saving and loading the file & & \bx \\ \midrule
Tests (\texttt{my\_tests.py} or \texttt{test\_mygame.py}) & & \bx \\ \midrule
Playtest notes and accessibility check & & \bx \\ \midrule
\rule{0pt}{1.2em} & & \bx \\
\bottomrule
\end{tabularx}}

\section{Planning page 4: My tests (M3)}
Test a \textbf{normal} case and at least one \textbf{edge} case (a wall, zero, an empty list, a
missing file).

{\renewcommand{\arraystretch}{1.8}
\begin{tabularx}{\linewidth}{@{}>{\raggedright\arraybackslash}p{3.6cm}|X|X|>{\centering\arraybackslash}p{1.9cm}@{}}
\toprule
\textbf{Function} & \textbf{I give it\ldots} & \textbf{It should return\ldots} & \textbf{Normal / edge} \\ \midrule
\texttt{clamp} (example) & \texttt{500, -270, 270} & \texttt{270} & edge \\ \midrule
 & & & \\ \midrule
 & & & \\ \midrule
 & & & \\
\bottomrule
\end{tabularx}}

\section{Planning page 5: Playtest and accessibility (M4)}
\begin{checklist}
  \item A playtester tried it while I watched quietly and took notes.
  \item I fixed or improved one thing from the playtest: \rule{4.5cm}{0.4pt}
  \item Colors come from the safe palette, and things that matter also differ in \textbf{shape} or \textbf{size}.
  \item Text is big enough to read from a chair.
  \item The keys are easy to reach, and the screen says which keys to use.
  \item It never asks for a real name, age, address, school or password.
  \item If I used someone else's idea, art or code, I wrote who made it in a comment.
\end{checklist}

%% ---------------------------------------------------------------------
\clearpage
\section{How your project is assessed}
Your teacher uses the same four questions. Check where you think you are now, and again at the
end.

{\small
\renewcommand{\arraystretch}{1.35}
\begin{longtable}{@{}>{\bfseries\raggedright\arraybackslash}p{2.3cm}>{\raggedright\arraybackslash}p{2.9cm}>{\raggedright\arraybackslash}p{2.9cm}>{\raggedright\arraybackslash}p{2.9cm}>{\raggedright\arraybackslash}p{3.0cm}@{}}
\toprule
 & Just starting & Getting there & Got it & Expert \\ \midrule\endhead
Creativity & I changed a few things in the template. & My project has my own theme and one new idea. & My project has its own idea, rules and look. & My idea has a surprising twist that makes people want another go. \\ \midrule
Working code & It runs, but crashes or the main part does not work yet. & The main part works; some must-haves are missing. & It works as my pitch said, with all six must-haves. & It also handles tricky moments (walls, missing file, playing again). \\ \midrule
Code quality and tests & My code is one long block; no tests yet. & I use some functions; 1--2 tests pass. & Small functions with clear names and comments; 3 or more tests pass. & Rules are separate from drawing; my tests include edge cases. \\ \midrule
Explaining it & I can show what it does. & With help, I can explain what one function does. & I can explain how one part works and one decision I made. & I can explain my design, what I would change, and a bug my tests caught. \\
\bottomrule
\end{longtable}}

\section{Support and extension}
\begin{differentiation}
\supportpath Start from the template and change it step by step. Aim first for M2 (the core
loop) and two tests; your teacher can be your navigator. A small project that works and that
you can explain is a great project.
\extensionpath Put your rules in their own file (like \texttt{starcatcher\_logic.py} in the
reference game) and import it. Add levels, a start menu, a top-5 high-score table, or let players
choose their own keys. Write a test for your file functions.
\end{differentiation}

\section{When you are stuck}
\begin{enumerate}[itemsep=2pt]
  \item Read the error message: start with the \textbf{last line}.
  \item Print the values: \py|print("x is", x)|.
  \item Explain your code line by line to your partner or a rubber duck.
  \item Make the problem smaller: test one function on its own.
  \item Check the cheat sheet and the \texttt{starter/README.md}.
  \item Ask your teacher: say what you expected and what happened instead.
\end{enumerate}

\begin{safety}[Showcase safety]
Use a nickname on screen and on your project card. Your project stays on our computers and is
shown to family only: nothing is uploaded. Credit anyone whose ideas, art or code you used.
\end{safety}

\end{document}
